\documentclass[11pt,letterpaper]{article}

\usepackage[T1]{fontenc}
\usepackage[utf8]{inputenc}
\usepackage{lmodern}
\usepackage{microtype}
\usepackage[margin=1in]{geometry}
\usepackage{booktabs}
\usepackage{tabularx}
\usepackage{longtable}
\usepackage{array}
\usepackage{enumitem}
\usepackage{listings}
\usepackage{xcolor}
\usepackage{url}
\usepackage[numbers]{natbib}
\usepackage[hidelinks,bookmarks=false]{hyperref}

\setlength{\parskip}{4pt}
\setlength{\parindent}{0pt}
\setlist[itemize]{leftmargin=*,nosep}

\definecolor{codebg}{gray}{0.96}
\lstset{
  basicstyle=\ttfamily\small,
  backgroundcolor=\color{codebg},
  frame=single,
  framerule=0pt,
  xleftmargin=1em,
  framexleftmargin=1em,
  breaklines=true,
  columns=fullflexible,
  keepspaces=true
}

\newcommand{\dagtoml}{DAG-TOML}
\newcommand{\measured}{\textsc{measured}}

\title{\bfseries A Stateful Executable Proof: Governing an HTTP Echo
Service Lifecycle Across Eight Languages}
\author{Werner Kasselman\\Verivus OSS\\\texttt{werner@verivus.com}}
\date{May 31, 2026}

\begin{document}
\maketitle

\begin{abstract}
A companion proof in this repository showed that a \dagtoml{} contract
can govern an instant, deterministic, stdout-only program across
languages. That contract could be checked by running each program to
completion and diffing stdout. This paper reports the follow-up: the
same \dagtoml{} machinery---contract declaration, implementation DAG,
traceability, readiness gate, evidence matrix, and executable
witnesses---pointed at a contract that \emph{cannot} be checked that
way. The contract is stated over eight languages, of which seven
(Go, Node/JavaScript, TypeScript, Python, C, Rust, Java) must each bind
a port, stay alive, echo a \texttt{POST} body byte-for-byte, answer
repeated requests, and shut down gracefully on
\texttt{SIGTERM}---finishing any in-flight request, releasing the port,
and exiting~0. The eighth language, AWK, is the declared signal boundary:
gawk exposes no clean script-level \texttt{SIGTERM} handler, so it is
held to the lifecycle contract only as a \textsc{skip}-with-rationale and
its echo is left unasserted. On the runner used for this paper the
load-bearing witness reported seven \textsc{pass} outcomes, one
\textsc{skip} for AWK, and zero \textsc{fail}. A negative-control witness confirms the graceful-shutdown
check is non-vacuous: it catches both a server that ignores
\texttt{SIGTERM} and a server that drops an in-flight response while
exiting~0. The contribution is not a benchmark or a production
framework. It is a compact, falsifiable artifact that names precisely
which parts of the proof structure had to grow for a process
\emph{lifecycle}---a lifecycle-aware witness harness, the \measured{}
result word used for real, and a port-as-shared-resource
discipline---and which dynamic guarantees the static validators still
cannot see. A pre-build runtime spike, reported in full, also corrects
one design assumption: Java's bundled \texttt{HttpServer} cannot set
\texttt{SO\_REUSEADDR}, which fixes the run order rather than the claim.
\end{abstract}

\section{Introduction}

Artifact-centered papers are easiest to review when claims, commands,
and files line up directly, and artifact-evaluation guidance makes the
same point operationally: reviewers need enough structure to connect a
paper's claims to executable scripts, source files, expected outputs,
and limitations~\cite{acmBadging2020,poplArtifact2023}. A companion
paper applied that structure to the smallest nontrivial instance: a
``Hello, world!'' contract requiring exact stdout bytes and exit~0.
That proof is honest but limited in one specific way---its contract is
satisfiable by a program that runs to completion, so the witness is a
single stdout \texttt{cmp}.

Real assurance targets are not like that. Services bind ports, stay
resident, answer over the network, and must shut down without dropping
work. This paper asks whether the same inspectable, falsifiable
\dagtoml{} proof structure survives that jump. The contract under study,
declared in \path{non-trivial-proof/proof-bundle}, is an HTTP echo
service. Each PASS-candidate implementation must bind
\texttt{127.0.0.1:8080}, return \texttt{200 OK} whose body is the exact
bytes of the request body, answer at least two sequential requests on one
process, and, on \texttt{SIGTERM}, finish any in-flight request, release
the port, and exit~0. AWK is exempted from this lifecycle contract as the
declared C06 boundary (Section~\ref{sec:observed} and the C06 row of
Table~\ref{tab:contracts}).

The jump forces machinery the stdout proof never needed. The verify
unit must background a server, poll for readiness, drive it over the
network, signal it, and capture the post-signal exit status with
guaranteed teardown. A bound port is shared mutable state on the runner,
so the language runs must be serialized and the port proven released
between them. Timing becomes a first-class observation, which is where
the fourth result word, \measured{}, is used for the first time in this
repository. And graceful shutdown is only meaningful against a negative
control, so the proof includes a witness whose \textsc{pass} means the
gate correctly classified a non-graceful server as \textsc{fail}.

This remains a minimal proof, not a performance study. There is no
throughput, latency percentile, or concurrency-under-load claim. The
goal is the same as before: an artifact a reviewer can inspect in one
sitting, with every dynamic guarantee tied to an executable witness and
every non-claim stated.

\section{Artifact}

The proof pack lives under \path{non-trivial-proof/proof-bundle}. Its
normative files are:

\begin{itemize}
  \item \path{contract_declaration.toml}: contracts C01 through C06.
  \item \path{implementation_dag.toml}: ten units, all tier 1.
  \item \path{traceability.toml}: three chains (service lifecycle,
        graceful-vs-kill, AWK boundary).
  \item \path{review_readiness.toml}: one gate for the service-witness pack.
  \item \path{evidence_matrix.toml}: five claims mapped to seven
        evidence artifacts.
  \item \path{run_service_contract.sh}, \path{detect_graceful_shutdown.sh},
        and \path{detect_awk_boundary.sh}: executable witnesses.
\end{itemize}

The eight servers are \path{src/go/server.go},
\path{src/node/server.js}, \path{src/typescript/server.ts},
\path{src/python/server.py}, \path{src/c/server.c},
\path{src/rust/server.rs}, \path{src/java/Server.java}, and
\path{src/awk/server.awk}. Two non-graceful negative controls live in
\path{src/controls/}. The echo payload is a fixed 56-byte JSON file,
\path{payload.json}.

\textbf{Eight languages, six runtimes, one boundary.} The firm headline
number is eight languages. ``Six runtimes'' is a derived figure under an
explicit definition: a runtime is a distinct language execution stack.
TypeScript folds into the V8/Node stack it shares with JavaScript, and
AWK is the C06 boundary rather than a passing runtime, leaving six
distinct stacks---Go, V8/Node (JS and TS), CPython, the JVM,
natively-compiled C, and natively-compiled Rust. We flag the soft edge
honestly: C and Rust are both bare native binaries over libc, so a
stricter ``managed-runtime'' definition would merge them and reach five.
The paper therefore leads with the number that is directly verifiable
from the source tree.

\section{Contract and Witnesses}

Table~\ref{tab:contracts} lists the six declared contracts and the local
witness that supports each. C01 is the load-bearing lifecycle contract.
C02 through C05 narrow separable obligations that depend on C01. C06 is
the declared AWK signal boundary with an explicit non-claim.

\begin{table}[htbp]
\centering
\small
\begin{tabularx}{\linewidth}{@{}p{0.10\linewidth}p{0.22\linewidth}X@{}}
\toprule
Contract & Domain & Witness and scope \\
\midrule
C01 & service\_lifecycle &
\path{run_service_contract.sh}; bind, byte-exact echo, exit 0 on
SIGTERM, port released. Load-bearing. \\
C02 & readiness &
Connectable on \texttt{127.0.0.1:8080} within 5000\,ms;
\emph{time-to-ready} reported as \measured{}. \\
C03 & echo\_fidelity &
Status exactly 200, body \texttt{cmp}-equal to the request body,
\texttt{Content-Length} equal to body length. \\
C04 & signal\_handling &
In-flight request completes, exit 0 within 3000\,ms, port re-bindable;
\emph{time-to-shutdown} \measured{}. Backed by the negative control. \\
C05 & statefulness &
One long-lived process answers at least two sequential requests. \\
C06 & signal\_boundary &
\path{detect_awk_boundary.sh}; gawk has no clean script-level SIGTERM
handler, so AWK C04 is a declared \textsc{skip}; its echo is
\textsc{unassessable}. \\
\bottomrule
\end{tabularx}
\caption{Contract-to-witness inventory.}
\label{tab:contracts}
\end{table}

The C04 in-flight check is made race-free by an explicit synchronization
point rather than a timing guess. The test-only \texttt{?delay\_ms=}
request path makes a handler flush its status line and
\texttt{Content-Length} headers immediately, then sleep before writing
the body. The witness blocks on reading those headers off the socket
(proof the handler has been entered and committed its response), and
only then sends \texttt{SIGTERM}. ``The request is in flight'' is thus
an observed fact, not an assumption about how fast the signal arrives.
The injected delay is fixed at 1000\,ms, strictly less than the
3000\,ms shutdown deadline, so a correct graceful server can both finish
the request and exit in time; the only way to fail is to actually drop
the in-flight response.

\section{Implementation DAG}

The implementation DAG is a fan-in graph, but here the fan-in is
load-bearing for correctness rather than tidiness: a single verifier
must own the shared port. Seven layer-0 build/prepare units fan into one
layer-1 verification unit that runs each server one at a time on
\texttt{:8080}. Two independent layer-0 leaves run the behavioural
sidecar witnesses.

\begin{lstlisting}
U01 go-server      \
U02 node/ts-server  \   (one Node/V8 unit, two source files)
U03 python-server    \
U04 c-server          > U08 verify-service-contract (Java first)
U05 rust-server      /
U06 java-server     /
U07 awk-server     /

U09 verify-graceful-vs-kill   (independent leaf)
U10 verify-awk-boundary       (independent leaf)
\end{lstlisting}

The computed values checked by the repository's implementation-DAG
validator are:

\begin{itemize}
  \item ten units;
  \item layer counts \texttt{\{0: 9, 1: 1\}};
  \item entry points \texttt{U01..U07, U09, U10} (empty \texttt{depends\_on});
  \item leaf nodes \texttt{U08, U09, U10} (empty \texttt{blocks});
        U09 and U10 are simultaneously entry points and leaves;
  \item critical path \texttt{U05 -> U08} with node-weighted
        critical-path LOC 327.
\end{itemize}

\textbf{Run order is not arbitrary.} U08 runs Java first, on a pristine
port. The reason is a measured runtime fact, not a preference; it is
discussed in Section~\ref{sec:java}.

\section{The Witness Harness (the part that grew)}

The stdout proof's harness ran each program to exit and compared bytes.
This harness adds five primitives the stdout proof never needed: PID
capture into a dedicated process group (so a failed assertion never
orphans a daemon or its children on the port), readiness polling with a
deadline, signal delivery, exit-after-signal capture, and guaranteed
teardown on every exit path. Port \texttt{8080} is treated as a global,
mutable, singleton resource: the verifier serializes the language runs,
fails fast to \textsc{skip} if the port is already occupied, and proves
the port is re-bindable after each server exits using an
\emph{independent} re-bind probe that itself sets \texttt{SO\_REUSEADDR}
and retries for up to two seconds to absorb \texttt{TIME\_WAIT}. Testing
only that the \emph{server} set the option is not enough, because the
probe is a fresh socket that can hit \texttt{TIME\_WAIT} on its own.

\section{Observed Execution}
\label{sec:observed}

All commands were run from the repository root on May 31, 2026. The full
toolchain was present: \texttt{go} 1.26.3, \texttt{node} v24.15.0,
\texttt{python3} 3.13.13, \texttt{cc} (gcc) 15.2.0, \texttt{rustc}
1.90.0, \texttt{java} 25.0.3 (runtime only; no \texttt{javac}), and
\texttt{gawk} 5.3.2, with port 8080 free.

\begin{table}[htbp]
\centering
\small
\begin{tabularx}{\linewidth}{@{}p{0.46\linewidth}p{0.07\linewidth}X@{}}
\toprule
Command & Exit & Observed result \\
\midrule
\path{bash .../run_service_contract.sh} & 0 &
7 pass, 1 skip, 0 fail. Go, Node, TS, Python, C, Rust, Java passed;
AWK skipped as the declared C06 boundary. \\
\path{bash .../detect_graceful_shutdown.sh} & 0 &
Both non-graceful controls caught as C04 FAIL (ignore-SIGTERM and
drop-in-flight). \\
\path{bash .../detect_awk_boundary.sh} & 0 &
gawk listener up; echo UNASSESSABLE; SIGTERM yields exit 143, confirming
no clean handler. \\
\path{validate_implementation_dag.py ... --check-paths-exist} & 0 &
Passed; 10 units, layers \texttt{\{0:9,1:1\}}, critical-path LOC 327. \\
\path{validate_traceability.py ... --check-paths-exist} & 0 &
Passed; 29 entities, path checks enabled. \\
\path{validate_review_readiness.py ...} (three files) & 0 &
Contract-declaration, readiness-gate, and evidence-matrix passed. \\
\bottomrule
\end{tabularx}
\caption{Observed command outcomes.}
\label{tab:commands}
\end{table}

Table~\ref{tab:lifecycle} gives the per-language lifecycle results from
one representative run. The two timing columns are \measured{} numbers,
not gates (Section~\ref{sec:measured}).

\begin{table}[htbp]
\centering
\small
\begin{tabularx}{\linewidth}{@{}lllX@{}}
\toprule
Language & Result & time-to-ready & time-to-shutdown \\
\midrule
Java   & PASS & 566\,ms & 24\,ms\,$\dagger$ \\
Go     & PASS & 4\,ms   & 1070\,ms \\
Node   & PASS & 48\,ms  & 1024\,ms \\
TS     & PASS & 71\,ms  & 1024\,ms \\
Python & PASS & 47\,ms  & 1026\,ms \\
C      & PASS & 3\,ms   & 1002\,ms \\
Rust   & PASS & 3\,ms   & 1001\,ms \\
AWK    & SKIP & ---     & C06 boundary; echo UNASSESSABLE \\
\bottomrule
\end{tabularx}
\caption{Per-language lifecycle (one representative run; \measured{}
timings vary run to run).}
\label{tab:lifecycle}
\end{table}

The roughly 1000\,ms shutdown for six runtimes is the injected 1000\,ms
in-flight delay being honoured: the server finishes the delayed request
before exiting. $\dagger$~Java is the exception at 24\,ms because
\texttt{server.stop(2)} interrupts the handler's sleep, so Java writes
the full body immediately instead of waiting out the artificial delay.
The in-flight response still arrives complete and byte-exact and the
exit code is 0, so C04 holds---this is honest graceful completion, in
direct contrast to the drop-in-flight control below.

\section{The \measured{} Result Word}
\label{sec:measured}

The companion stdout proof used four result words but reported no
\measured{} values, because nothing in that contract was a timed
observation. This proof is the first in the repository to use the fourth
word for real, and it is disciplined about what the word means. A
\measured{} number---\emph{time-to-ready}, \emph{time-to-shutdown}, or
bytes echoed---is descriptive evidence, never a \textsc{pass} on its
own. ``Within the deadline'' is a \emph{separate} \textsc{pass}/\textsc{fail}
contract: C02 fails if the server is not connectable within 5000\,ms,
and C04 fails if it has not exited within 3000\,ms, regardless of the
exact measured figure. The timings in Table~\ref{tab:lifecycle} vary
from run to run; that variance is precisely why they are reported as
measurements and not promoted to gates. (A scope note: the claim of
``first \measured{} use'' is currently scoped to this papers repository,
not the wider organization.)

\section{Negative Control: Graceful versus Killed}

A graceful-shutdown check that only inspects the exit code is vacuous: a
server can \texttt{exit(0)} the instant it receives \texttt{SIGTERM} and
look clean while dropping every in-flight request. The C04 definition is
therefore the stronger one---an in-flight request must \emph{complete}
---and \path{detect_graceful_shutdown.sh} proves the gate enforces it on
both axes. It drives two deliberately broken Go servers:

\begin{itemize}
  \item \textbf{Control A} installs no \texttt{SIGTERM} handler. The
        witness must find it still running after the deadline and
        escalate to \texttt{SIGKILL}; that is a C04 \textsc{fail}.
  \item \textbf{Control B} calls \texttt{exit(0)} on \texttt{SIGTERM}
        while a slow request is mid-flight. Its exit code is a clean 0,
        but the client receives a truncated body. The witness must read
        the truncated response and classify it \textsc{fail} anyway.
\end{itemize}

On this runner both controls were caught: Control A timed out and needed
\texttt{SIGKILL}; Control B exited 0 but delivered 0 of 56 expected body
bytes. \textsc{pass} for this witness means the C04 check discriminates
graceful from non-graceful on the exit-code axis \emph{and} the
in-flight-completion axis. This is the single most important honesty
witness in the bundle, and it is the lifecycle analogue of the stdout
proof's ``a divergent implementation would fail'' control.

\section{A Measured Runtime Correction: Java and SO\_REUSEADDR}
\label{sec:java}

Service proofs interact with the operating system in ways a static
design document can get wrong, so the design was checked against a
pre-build runtime spike before any server was written. The spike
overturned one design assumption and is reported here in full because
correcting it is part of the result.

The draft design assumed Java's bundled
\texttt{com.sun.net.httpserver.HttpServer} would carry the same
\texttt{SO\_REUSEADDR} guarantee as the other runtimes. It does not. On
this runner (OpenJDK 25.0.3, Linux, \texttt{tcp\_fin\_timeout}=60) the
spike established that \texttt{HttpServer} sets no \texttt{SO\_REUSEADDR}
and exposes no API to set it: neither \texttt{HttpServer.create(addr,
backlog)} nor the unbound \texttt{create()} followed by \texttt{bind()}
surfaces the underlying \texttt{ServerSocket}, so
\texttt{setReuseAddress(true)} before \texttt{bind} is unreachable. Once
the server has served a connection and closed it server-side, the port
holds a \texttt{TIME\_WAIT} entry and a fresh \texttt{HttpServer} cannot
re-bind it for the full timeout window; every spike run ended in
\texttt{BindException: Address already in use}. A plain
\texttt{ServerSocket} with \texttt{setReuseAddress(true)} binds the same
port immediately.

The correction is a run-order constraint, not a weakened claim. Java is
scheduled \emph{first} in the serialized verify order, on a pristine
port; its pre-flight is a \texttt{TIME\_WAIT}-aware probe rather than the
connectivity probe used for the other six. A lingering foreign
\texttt{TIME\_WAIT} when Java's turn comes is recorded as a legitimate
\textsc{skip}, never a false \textsc{fail}. Java's own port-release check
still passes, because that assertion is made by the independent
\texttt{SO\_REUSEADDR}-setting probe socket, not by Java re-binding. The
other six runtimes set \texttt{SO\_REUSEADDR} (Go and Node by default, C
via \texttt{setsockopt}, Rust via a small \texttt{extern "C"} FFI block,
Python via \texttt{allow\_reuse\_address}) and tolerate a prior
\texttt{TIME\_WAIT}.

\section{PASS, SKIP, FAIL, and MEASURED}

The proof uses the four result words narrowly.

\begin{itemize}
  \item \textsc{pass} means the witness exited 0 and reported the checked
        lifecycle condition held: ready, byte-exact echo over at least
        two requests, in-flight completion, exit 0 within the deadline,
        and port released.
  \item \textsc{skip} means a declared check was not performed because a
        toolchain was absent, the port was occupied, or the runtime is
        the declared C06 boundary (AWK). \textsc{skip} is not evidence of
        conformance.
  \item \textsc{fail} means a bind failure, a missed deadline, a non-200
        status, a body mismatch, a dropped in-flight response, a nonzero
        exit on \texttt{SIGTERM}, or a port not released.
  \item \measured{} is reserved for the timing and byte observations of
        Section~\ref{sec:measured}; it is never promoted to a
        \textsc{pass}.
\end{itemize}

\section{Claim Audit}

Table~\ref{tab:claims} maps the paper's claims to evidence, separating
direct observation from inference.

\begin{longtable}{@{}p{0.26\linewidth}p{0.28\linewidth}p{0.18\linewidth}p{0.18\linewidth}@{}}
\toprule
Claim & Evidence source & Status & Counterexample boundary \\
\midrule
\endhead
The proof pack is structurally complete for this example. &
the five \path{*.toml} files and the repository validators. &
Directly observed. &
Does not imply production readiness. \\
The seven available servers satisfy the C01--C05 lifecycle on this
runner. &
\path{run_service_contract.sh} and the observed run. &
Directly observed; AWK skipped as the C06 boundary. &
A runner missing a toolchain or port would produce more SKIPs. \\
The C04 graceful check is non-vacuous. &
\path{detect_graceful_shutdown.sh} and the two controls. &
Directly observed. &
Does not prove zero drops under arbitrary concurrency. \\
Java's \texttt{HttpServer} cannot set \texttt{SO\_REUSEADDR}; Java is
scheduled first. &
the pre-build spike and the witness pre-flight. &
Directly observed on this JDK and runner. &
Does not characterize other JDKs or HTTP libraries. \\
The AWK signal boundary is real and its echo is unasserted. &
\path{detect_awk_boundary.sh}; SIGTERM exit 143. &
Directly observed plus a declared non-claim. &
Does not prove gawk \texttt{/inet} cannot echo byte-exact. \\
Artifact organization follows current review expectations. &
ACM badging, POPL artifact guidance, and SIGSOFT
standards~\cite{acmBadging2020,poplArtifact2023,sigsoftStandards}. &
Cited. &
No formal artifact badge is claimed. \\
\bottomrule
\caption{Claim-to-evidence audit.}
\label{tab:claims}
\end{longtable}

\section{Related Work}

The proof's framing is an executable specification used as a consistency
check rather than a benchmark, a line of work with precedent in trace
specifications, where formal specifications are turned into executable
models~\cite{hoffman1988trace}. The lifecycle contract itself encodes
widely-recommended operational practice: disposable processes that start
fast and shut down gracefully on \texttt{SIGTERM}, finishing in-flight
work before exit, is the disposability principle of the twelve-factor
methodology~\cite{twelvefactor}, and treating non-graceful termination
as a first-class failure mode to be tested echoes stability-pattern
guidance for production services~\cite{nygard2018release}.

The artifact framing follows research-artifact guidance rather than
benchmark methodology. ACM distinguishes artifact evaluation,
availability, and result validation~\cite{acmBadging2020}; POPL's
artifact guidance asks authors to map scripts and source files to paper
claims~\cite{poplArtifact2023}; and SIGSOFT's empirical standards
emphasize specific, technical review checklists for software-engineering
research~\cite{sigsoftStandards,ralph2020empirical}.

\section{Threats to Validity}

The validity structure separates what was observed from the broader
claims one might draw~\cite{runeson2009case,feldt2010validity}.

\textbf{Construct validity.} The proof measures whether a stateful
lifecycle contract is represented and checked, not whether \dagtoml{} is
sufficient for a large assurance pipeline. The echo service is
intentionally minimal: it is not a benchmark, not RFC-conformant HTTP,
and has no TLS or keep-alive guarantees beyond the test.

\textbf{Internal validity.} The strongest internal risks were addressed
before the reported run. The in-flight check uses an explicit
header-flush synchronization point rather than a timing guess; the
port-release check uses an independent \texttt{SO\_REUSEADDR} re-bind
probe rather than a mere listener-absence check; and the byte-exact
comparison checks \texttt{Content-Length} as well as the body. Each of
these was a real defect caught and corrected during review.

\textbf{External validity.} The proof covers eight small servers on one
POSIX runner. \texttt{SIGTERM}/\texttt{SIGKILL} semantics assume POSIX;
Windows is out of scope. Port 8080 is a runner-global assumption, and a
hostile environment racing for the port yields a legitimate \textsc{skip}.

\textbf{Conclusion validity.} The results are categorical witness
outcomes and descriptive \measured{} timings, not statistical estimates.
The correct conclusion is that the declared witnesses passed, skipped, or
failed as reported on this runner.

\section{What the Static Validators Do Not Check}

The five \dagtoml{} validators check structure, not runtime. They
confirm that C01 is declared, wired to a test, and that the witness
scripts exist; they cannot confirm that a server ever bound, echoed, or
shut down. Every dynamic guarantee lives only in the shell witnesses and
is surfaced into Section~\ref{sec:observed}. The readiness-gate
validator, in particular, requires a gate identifier, an artifact class,
and one of a small set of content fields, and resolves the artifact-class
link and the required-document paths; it does \emph{not} interpret the
human-readable pass and block conditions, which are enforced only by
running the witness. Path existence is itself opt-in: every validator
invocation in this paper passes \texttt{--check-paths-exist} with a
repository root, or the existence guarantee would be hollow. There is no
validator hook for \measured{}; it is a paper-level reporting discipline.
This gap is the natural bridge to later follow-ups in the same arc,
especially handing verification to enterprise static-analysis tools, and
is noted here as future work without overclaiming.

\section{arXiv and Artifact Packaging Notes}

The public artifact repository for this paper is
\url{https://github.com/verivus-oss/agent-assurance-papers}; the
specification and validators are at
\url{https://github.com/verivus-oss/agent-assurance}. The package is
intentionally plain: \texttt{article}, \texttt{pdflatex}, BibTeX,
\texttt{natbib}, and standard packages. arXiv processes TeX submissions
automatically, expects authors to inspect the generated PDF, and asks
that extraneous files such as logs and aux files be removed from the
source package~\cite{arxivSubmitTex2026,arxivTexLive2026}. An arXiv
source package for this paper should include only \path{main.tex},
\path{references.bib}, and a precomputed \path{main.bbl}; the bundle code
can be linked or attached as ancillary material.

\section{Limitations and Non-Claims}

This paper does not claim throughput, latency, or concurrency-under-load
results; RFC-conformant HTTP, TLS, or keep-alive guarantees; correctness
on non-POSIX runners; that gawk \texttt{/inet} cannot achieve byte-exact
echo (only that this bundle does not assert it); or production readiness.
The exact claim is narrower: the repository contains a multi-language
stateful-service proof whose lifecycle contract, implementation graph,
traceability chains, and review gate can be checked against executable
witnesses; whose graceful-shutdown check is shown non-vacuous by negative
controls; whose one declared signal boundary is explicit; and whose
observed outcomes are reproducible by running the listed commands in an
environment with the same toolchains.

\section{Conclusion}

The stdout proof showed that a \dagtoml{} contract can govern identical
output bytes across languages. This proof shows the same structure
governing an identical \emph{lifecycle}---ready, echo, graceful
shutdown, port release---across seven of eight languages (six runtimes),
with the eighth, AWK, held as the one declared signal boundary rather
than to the lifecycle contract. What had to grow is named precisely: a
lifecycle-aware witness harness, the \measured{} result word used for
real, a port-as-shared-resource discipline, and a negative control that
keeps the graceful-shutdown check honest. What did not change is the
proof structure itself. The artifact is not evidence about performance or
production assurance; it is a compact, executable proof that a stateful
service contract can be tied to contracts, DAG nodes, code paths,
witnesses, and review gates, and that the parts a static validator cannot
see are stated rather than hidden.

\bibliographystyle{plainnat}
\bibliography{references}

\end{document}

\typeout{get arXiv to do 4 passes: Label(s) may have changed. Rerun to get cross-references right.}
